\documentclass[../thesis.tex]{subfiles}

\begin{document}
	Здесь будет детально рассмотрена реализация всех подсистем и их взаимодействие. В соответствии с
	критериями сравнения будут выбраны подходящие технологии. На рисунке \ref{fig:3} схематично изображена
	структура системы.

	\begin{figure}[H]
		\centering
		\includegraphics[width=1\textwidth]{assets/img/p2/impl.png}
		\caption{Структура разработанной системы}
		\label{fig:3}
	\end{figure}

	\section{Менеджер сервисов}

	Тут будет детальное описание модели работы с гетерогенными сервисами. С техническим описанием
	менеджера сервисов можно ознакомится в Приложении \hyperref[apx:1]{А}.

	\section{Графический интерфейс}

	\subsection{Платформа оператора}

	В качестве платформы операторы используется Appsmith. Для разработки приложений предоставляется конструктор
	с большим выбором готовых графических компонентов.

	Взаимодействие с менеджером сервисов осуществляется через REST API. Для этих целей в Appsmith
	предусмотрены возможности для написания запросов на JavaScript. Чтобы упростить процесс разработки
	и сократить размер дублирующегося функционала был разработан универсальный набор JavaScript объектов,
	которые инкапсулируют взаимодействие с сервером.

	\subsection{Визуализация}

	Для визуализации результатов работы сервиса (численных вычислений и информационных сообщений) используется
	Grafana, предоставляющую широкие возможности для интерпретации научных экспериментов (см.
	\cite{grafana-cern}). Для каждого сервиса разработчик конструирует панель, которая интегрируется в
	платформу оператора (приложение Appsmith) в виде iFrame.

	На момент создания системы, в единственном открытом расширении
	\linebreak
	Grafana (см. \cite{grafana-cern}) с поддержкой Apache Kafka не было поддержки работы с
	\linebreak
	Protobuf. В связи с этим в расширение был интегрирован функционал Redpanda Connect для
	построения потоков данных с использование регистра схем. Вместо описания в \textit{YAML} файле, поток
	с заданными параметрами программно конструируется внутри расширения.
\end{document}